\subsection*{一步一步建立直观}

\paragraph{1. 先把向量看成一次位移}

想象在一张地图上先走 $u$，再走 $v$。第二段路可以搬到第一段的终点，方向和长度却不改变。右图的虚线正是在做这个搬移；红色箭头只记录起点到最终位置的总变化。它不记录走过的路径，所以同一位移可以由不同路线实现。



\paragraph{2. 数乘改变的是位移的份数}

把系数 $s$ 改为二，相当于沿 $u$ 连走两次，再走一次 $v$。总位移是 $su+v$，横纵两个分量分别累加。这里的加法并非任意规定：它来自每个坐标方向上净位移的相加，因此与先后分组无关。负系数则把相应位移反向。



\paragraph{3. 方向数目与独立信息}

让两根箭头共线后，虽然仍可以填写两个系数，终点却只能沿一条直线移动。图形帮助我们分辨两个数字和两份独立信息。若两根箭头不共线，改变两个系数才能分别控制平面内的两个方向；这就是下一节讨论基的原因。



\subsection*{定义、命题与证明}

\paragraph{定义：线性组合}

给定向量 $u,v$ 与标量 $s,t$，向量 $su+tv$ 称为它们的线性组合。系数记录每个方向参与的份额；负系数表示反向。

\paragraph{命题：向量平移不改变合位移}

在同一向量空间中，先走 $u$ 再走 $v$ 的总位移为 $u+v$，与先走 $v$ 再走 $u$ 相同。

\begin{proof}

用坐标表示 $u=(u_1,u_2)$、$v=(v_1,v_2)$。第一条路径的终点是 $(u_1+v_1,u_2+v_2)$，第二条路径的终点是 $(v_1+u_1,v_2+u_2)$。

实数加法交换律使两点一致。这里比较的是终点位移；两条路径的中间位置可以不同，所以不能把结论改写为两次运动在每一时刻都相同。

\end{proof}
